\section{Hai quy luật, và chúng không đồng ý}
\label{sec:ch09-hai-quy-luat}

\index{quy luật scaling!ba luật lũy thừa của Kaplan}%
Nếu phát kiến là \tn{tiền huấn luyện}{pretraining} trên văn bản không nhãn, thì
câu hỏi tiếp theo là câu hỏi kỹ thuật: có một ngân sách tính toán, tiêu vào đâu.

Hai đáp án, cách nhau hai năm, và chúng ngược nhau.

\subsection*{Kaplan 2020: mất mát là hàm lũy thừa của ba thứ}

Bài Kaplan khớp ba luật, mỗi luật cho một trong ba đại lượng, và mục 1.2 in cả
ba với dấu \(\sim\) chứ không phải dấu bằng:

\begin{align}
\label{eq:ch09-kaplan}
L(N) &= (N_c/N)^{\alpha_N}, & \alpha_N &\sim 0.076, & N_c &\sim 8.8 \times 10^{13}, \\
L(D) &= (D_c/D)^{\alpha_D}, & \alpha_D &\sim 0.095, & D_c &\sim 5.4 \times 10^{13}, \notag \\
L(C_{\min}) &= (C_c/C_{\min})^{\alpha_C}, & \alpha_C &\sim 0.050, & C_c &\sim 3.1 \times 10^{8}. \notag
\end{align}

\(N\) là tham số không tính embedding, \(D\) là token, \(C_{\min}\) tính bằng
petaflop-ngày. Hai chữ số có nghĩa là tất cả những gì bài in, và tôi không thêm
chữ số thứ ba.

Số mũ nhỏ tới mức đáng ngờ. \(\alpha_N = 0.076\) nghĩa là tăng mô hình gấp mười
thì mất mát chỉ nhân với \(10^{-0.076}\), tức giảm 16.1\%. Đó chính là
hình dạng của mọi thứ trong bài: cải thiện có, đều đặn, và đắt.

\begin{measured}{ba luật của Kaplan, tính ra số}
\begin{minted}[fontsize=\footnotesize]{text}
N (non-embedding)   L(N)      D (tokens)      L(D)      C_min (PF-days)  L(C)
1.0e+07             3.3712    1.0e+08         3.5041    1.0e+00          2.6581
1.0e+08             2.8300    1.0e+09         2.8156    1.0e+02          2.1114
1.0e+09             2.3756    1.0e+10         2.2624    1.0e+04          1.6771
1.0e+10             1.9943    1.0e+11         1.8179    1.0e+06          1.3322
\end{minted}

Đo bằng \code{experiments/ch09\_laws.py} tại tag \code{ch09}.
\end{measured}

\textbf{Một cái bẫy khi trích bài này.} Cùng những ký hiệu ấy mang giá trị khác
nhau ở những chỗ khác nhau trong bài. Bảng 3 mục 5 khớp dạng hai biến
\(L(N, S)\) và cho \(\alpha_N = 0.077\) với \(N_c = 6.5 \times 10^{13}\); phụ
lục A bảng 5 khớp \(L(C)\) dạng thô và cho \(\alpha_C = 0.057\). Bài không mâu
thuẫn với chính nó, vì đó là ba phép khớp cho ba mục đích, và chú thích 3 nói
rõ nên dùng đường \(C_{\min}\) để dự đoán. Nhưng trích \(N_c\) mà không nói lấy
từ phương trình nào là trích sai, và đây là dạng của quyết định về việc trích
theo phiên bản mà chương~\ref{ch:transformer} đã phải học một lần.

Câu trả lời cho câu hỏi ngân sách nằm ở bảng 6 phụ lục A: \(N_{\mathrm{opt}}
\propto C_{\min}^{0.73}\) và \(D_{\mathrm{opt}} \propto C_{\min}^{0.27}\). Bài
nói bằng lời ở trang 5: \enquote{As the computational budget \(C\) increases, it
should be spent primarily on larger models, without dramatic increases in
training time or dataset size.}

Tiêu vào tham số. Đó là lời khuyên, và ngành nghe theo: GPT-3 ra cùng năm ấy,
với 175 tỷ tham số trên 300 tỷ token.

\subsection*{Hoffmann 2022: chia đôi}

\index{quy luật scaling!Chinchilla và phép chia đôi}%
Hai năm sau, cùng câu hỏi, 400 mô hình. Tóm tắt:

\begin{quote}
\enquote{we find that for compute-optimal training, the model size and the
number of training tokens should be scaled equally: for every doubling of model
size the number of training tokens should also be
doubled.}~\autocite{hoffmann2022chinchilla}
\end{quote}

Bảng 2 của bài cho ba cách ước lượng và ba cặp số mũ: 0.50 và 0.50 cho cách
thứ nhất, 0.49 và 0.51 cho cách thứ hai, 0.46 và 0.54 cho cách thứ ba, với phân
vị bootstrap thứ 10 và 90 in kèm. Hàng cuối bảng là chính Kaplan, 0.73 và 0.27.

Nói bằng số mũ thì hai bên nghe như chỉ khác nhau về mức độ. Nói bằng số nhân
thì khác hẳn.

\begin{measured}{một hệ số ngân sách đi về đâu, dưới mỗi bài}
\begin{minted}{text}
compute x   Kaplan: model x   Kaplan: data x   Hoffmann: each x
2           1.66              1.21             1.41
10          5.37              1.86             3.16
100         28.84             3.47             10.00
1000        154.88            6.46             31.62
\end{minted}

Đo bằng \code{experiments/ch09\_laws.py} tại tag \code{ch09}.
\end{measured}

Ngân sách gấp nghìn: Kaplan nói mô hình lớn gấp 155 lần với dữ liệu chỉ gấp 6.5;
Hoffmann nói cả hai gấp 31.6. Ở hàng cuối ấy hai lời khuyên cho ra hai mô hình
lệch nhau năm lần về kích thước.

Một chi tiết nhỏ trong bảng trên. Phần mở đầu của bài Hoffmann diễn đạt lại
hàng \code{10} của Kaplan thành \enquote{5.5x} và \enquote{1.8x}. Số mũ của
chính Kaplan cho 5.37 và 1.86. Bài làm tròn số của đối thủ chứ không trích nó,
và chênh lệch không đổi kết luận nào; tôi ghi lại vì một người đọc dựng lại
bảng trên sẽ ra 5.37 và tự hỏi mình sai ở đâu.

\subsection*{Một chỗ bài không dựng lại được từ chính nó}

\index{quy luật scaling!hằng số làm tròn của Hoffmann}%
Hàng thứ ba của bảng 2 không phải một phép khớp độc lập. Phương trình 4 dẫn nó
ra từ \tn{hàm mất mát}{loss function} hai biến của phụ lục D.2,

\begin{equation}
\label{eq:ch09-chinchilla}
L(N, D) = E + \frac{A}{N^{\alpha}} + \frac{B}{D^{\beta}},
\end{equation}

với \(E = 1.69\), \(A = 406.4\), \(B = 410.7\), \(\alpha = 0.34\) và
\(\beta = 0.28\). Từ đó \(a = \beta/(\alpha + \beta)\) và
\(b = \alpha/(\alpha + \beta)\), nên \(a + b = 1\) đúng theo cấu trúc, và đó
mới là phát biểu thật: một hệ số ngân sách chia hết cho mô hình và dữ liệu,
không đi đâu khác.

Thay số vào thì \(a = 0.4516\) và \(b = 0.5484\). Bảng 2 in 0.46 và 0.54, và
0.4516 không làm tròn thành 0.46.

\begin{measured}{giải ngược ra \(\alpha\) và \(\beta\) mà bảng 2 hàm ý}
\begin{minted}{text}
  a = beta/(alpha+beta)  = 0.4516  rounds to 0.45
  b = alpha/(alpha+beta) = 0.5484  rounds to 0.55
  a + b = 1.0000  (exactly 1 by construction)

  beta/alpha must be          0.8519
  with alpha = 0.34, beta would be  0.2896
  with beta = 0.28, alpha would be   0.3287
\end{minted}

Đo bằng \code{experiments/ch09\_laws.py} tại tag \code{ch09}.
\end{measured}

Giữ \(\alpha\) đã in thì \(\beta\) phải là 0.2896; giữ \(\beta\) đã in thì
\(\alpha\) phải là 0.3287. Không giá trị nào sống sót một mình, nên cả hai đều
mang thêm chữ số mà phụ lục D.2 không in, và đường biên tính trước khi làm tròn.

Đây không phải mâu thuẫn, và tôi không đọc nó thành lỗi. Nó là chuyện khác:
\eqref{eq:ch09-chinchilla} và bảng 2 không cùng dựng lại được từ hằng số bài in
ra. Bạn tính được hàm mất mát, hoặc tính được đường biên, không phải cả hai.

\subsection*{Chinchilla so với Gopher}

\index{quy luật scaling!Chinchilla so với Gopher}%
Bài không dừng ở phép khớp. Nó lấy ngân sách của một mô hình đã có và tiêu lại
theo công thức mới. Gopher là 280 tỷ tham số trên 300 tỷ token; Chinchilla là 70
tỷ trên 1.4 nghìn tỷ. Bốn lần nhỏ hơn, gần năm lần nhiều dữ liệu hơn.

Kết quả trên Massive Multitask Language Understanding, viết tắt là MMLU: Gopher
60.0, Chinchilla 67.6, theo bảng 6 mục 4.2.2. Tóm tắt
của bài in 67.5 và \enquote{greater than a 7\% improvement}, còn thân bài in
67.6 và \enquote{improving upon Gopher by 7.6\%}. Hai con số ấy lệch nhau trong
chính bài, giống hệt ở cả bản arXiv lẫn bản NeurIPS. Chương này dùng 67.6, vì đó
là số bảng 6 đo được, và ghi lại chỗ lệch thay vì chọn im lặng.

\begin{measured}{cùng một ngân sách, chia hai kiểu}
\begin{minted}{text}
model         N            D            C = 6ND      L(N, D)
Gopher        2.8e+11      3e+11        5.040e+23    1.9933
Chinchilla    7e+10        1.4e+12      5.880e+23    1.9366
\end{minted}

\eqref{eq:ch09-chinchilla} tính ở hai cấu hình bảng 1 của bài in.

Đo bằng \code{experiments/ch09\_laws.py} tại tag \code{ch09}.
\end{measured}

Hàm mất mát của chính bài xếp Chinchilla trên Gopher, đúng chiều bài đo được
trên MMLU. Nhưng cột \code{C} thì không khớp với điều bài nói về nó. Mục 4 viết
hai mô hình \enquote{have been trained for the same number of FLOPs}, và
\(6ND\) trên chính con số bảng 1 cho 5.040e+23 với 5.880e+23, lệch nhau 16.7\%.
Để bằng nhau dưới \(6ND\), Chinchilla cần 1.200e+12 token thay vì 1.4e+12.

Cách giải thích khả dĩ nằm trong chính bài, và nó cũng là lý do tôi không đọc
chỗ này thành lỗi. Bài không dùng \(6ND\) cho phép đếm của mình: phụ lục F đếm
đủ mọi số hạng, kể cả embedding và logits cuối, và bảng A4 của bài đo tỷ lệ
giữa phép đếm đủ ấy với \(6ND\), rồi báo rằng nó chạy từ 0.99 tới 1.10 tùy hình
dạng mô hình. Gopher rộng gấp đôi Chinchilla, nên hai mô hình nằm ở hai chỗ khác
nhau của dải ấy, và một khoảng 16.7\% dưới \(6ND\) không đòi hỏi con số nào
phải sai.

\textbf{Chỗ này tôi định kiểm và không kiểm được.} Bảng A4 không dựng lại được
từ bài, vì phép đếm của phụ lục F phụ thuộc cả độ dài chuỗi lẫn kích thước từ
điển, và bài không in con số nào trong hai. Tôi tìm cả hai trong mục huấn luyện
và trong từng phụ lục từ A tới I, và không thấy.

Nên tôi làm việc còn lại làm được: cài phép đếm của phụ lục F rồi quét hai ẩn
ấy trên chín cặp giá trị hợp lý, ba độ dài chuỗi nhân ba kích thước từ điển.

\begin{measured}{chín cặp, và không cặp nào cho lại cột của bảng A4}
\begin{minted}{text}
published:    1.03   1.10   1.08   1.04   1.03   0.99

n_ctx  vocab   rebuilt ratios                          worst gap
1024   32000    1.48   1.26   1.20   1.14   1.12   1.07   0.455
2048   32000    1.68   1.41   1.33   1.24   1.20   1.11   0.645
4096   32000    2.06   1.71   1.57   1.42   1.36   1.21   1.026
\end{minted}

Ba hàng đại diện của chín; \code{ch09\_laws.py} in cả chín. Cặp gần nhất vẫn
lệch 0.455 ở hàng tệ nhất của nó.

Đo bằng \code{experiments/ch09\_laws.py} tại tag \code{ch09}.
\end{measured}

Chín cặp, không cặp nào tới gần. Nên chương này ghi lại bảng A4 nói điều gì, và
ghi lại rằng nó không kiểm được điều đó, thay vì in một cột số của riêng mình
cạnh cột của bài như thể hai cột so được với nhau.

\subsection*{GPT-3 dưới thước đo của Chinchilla}

\index{quy luật scaling!GPT-3 huấn luyện thiếu}%
Bảng 3 của bài chiếu công thức mới lên chín ngân sách. Hàng 175 tỷ tham số muốn
3.7 nghìn tỷ token. GPT-3 dùng 300 tỷ, kém 12.3 lần.

\begin{measured}{GPT-3 và một mô hình cùng ngân sách, dưới \eqref{eq:ch09-chinchilla}}
\begin{minted}{text}
model                      N            D            L(N, D)
GPT-3 as trained           1.75e+11     3e+11        2.0023
compute-matched, 20 D/N    5.123e+10    1.025e+12    1.9609
\end{minted}

Mô hình đối chứng dựng từ chính tỷ số của bảng 3: \(C = 6N(20N) = 120N^2\), nên
\(N = \sqrt{C/120}\), ở đúng \(C = 3.150 \times 10^{23}\) của GPT-3.

Đo bằng \code{experiments/ch09\_laws.py} tại tag \code{ch09}.
\end{measured}

Cùng ngân sách, và phép khớp thích mô hình 51 tỷ tham số hơn mô hình 175 tỷ,
hơn 0.0414 nat. Đọc theo chiều thực hành: theo bài này, cái đắt nhất GPT-3 mua
là kích thước nó không dùng hết.

Tôi để một cảnh báo cạnh đoạn ấy. \eqref{eq:ch09-chinchilla} là một phép khớp
trên dữ liệu của một phòng thí nghiệm, với tokenizer của họ và corpus của họ, và
GPT-3 huấn luyện trên corpus khác với tokenizer khác. Con số 0.0414 nat là con
số của mô hình đã khớp, không phải một phép đo trên GPT-3. Đó là lý do nó nằm
trong một hộp \enquote{tôi đo được} ghi rõ công thức, chứ không nằm trong một
câu bảo rằng GPT-3 lẽ ra nên nhỏ hơn.

\subsection*{Quy tắc hai mươi token, và chỗ nó không được phát biểu}

\index{quy luật scaling!quy tắc hai mươi token}%
Bài này hay được rút gọn thành một quy tắc: chừng hai mươi token cho mỗi tham
số. Quy tắc ấy tiện, và nó không có trong bài.

\begin{measured}{bảng 3 của bài, cộng hai cột tôi tính thêm}
\begin{minted}[fontsize=\footnotesize]{text}
parameters     tokens        D/N     6ND          table's FLOPs  ratio
400 Million    8.000e+09     20.0    1.920e+19    1.920e+19      1.0000
1 Billion      2.020e+10     20.2    1.212e+20    1.210e+20      1.0017
10 Billion     2.051e+11     20.5    1.231e+22    1.230e+22      1.0005
67 Billion     1.500e+12     22.4    6.030e+23    5.760e+23      1.0469
175 Billion    3.700e+12     21.1    3.885e+24    3.850e+24      1.0091
280 Billion    5.900e+12     21.1    9.912e+24    9.900e+24      1.0012
520 Billion    1.100e+13     21.2    3.432e+25    3.430e+25      1.0006
1 Trillion     2.120e+13     21.2    1.272e+26    1.270e+26      1.0016
10 Trillion    2.162e+14     21.6    1.297e+28    1.300e+28      0.9978
\end{minted}

Đo bằng \code{experiments/ch09\_laws.py} tại tag \code{ch09}.
\end{measured}

Cột \code{D/N} chạy từ 20.0 tới 22.4 và trôi lên theo quy mô, chứ không đứng
yên ở một con số. Tôi tìm trong phần chính và trong từng phụ lục A tới I và
không thấy bài
viết quy tắc ấy ra ở đâu; nó là một cách đọc cột này. Chương gọi đúng tên như
vậy, theo cùng nguyên tắc mà mục~\ref{sec:ch07-chia-can} đã phải giữ một lần:
một câu chưa được phát biểu thì không được trích như thể đã.

Cột cuối cùng là một phép kiểm phụ và nó trả lời được một câu: bảng 3 dựng bằng
\(6ND\), khớp tới nửa phần trăm ở bảy trong chín hàng. Hai hàng còn lại lệch vì
cột token của chúng làm tròn: hàng 67 tỷ lệch 4.7\% ở 1.5 nghìn tỷ token, và
hàng 175 tỷ lệch 0.91\% ở 3.7 nghìn tỷ.

\subsection*{Vì sao hai bài lệch nhau, theo lời bài sau}

\index{quy luật scaling!lý do hai bài lệch nhau}%
Bài Hoffmann giải thích chỗ lệch, ở mục 2, và đưa đúng hai lý do.

Lý do thứ nhất là lịch learning rate:

\begin{quote}
\enquote{First, the authors use a fixed number of training tokens and learning
rate schedule for all models; this prevents them from modelling the impact of
these hyperparameters on the loss. [...] For a fixed learning rate cosine
schedule to 130B tokens, the intermediate loss estimates (for \(D' \ll\) 130B)
are therefore overestimates of the loss of a model trained with a schedule
length matching \(D'\).}
\end{quote}

Đọc chậm thì cơ chế rõ. Kaplan chạy mọi mô hình với một lịch cosine đặt cho một
độ dài cố định, và lấy mất mát ở các mốc giữa chừng làm điểm dữ liệu cho những
độ dài huấn luyện ngắn hơn. Nhưng một mô hình dừng ở giữa một lịch dài chưa
được hạ learning rate xuống đáy, nên nó tệ hơn một mô hình có lịch riêng vừa
đúng độ dài ấy. Mọi điểm dữ liệu \enquote{ít dữ liệu} vì thế đều bị đánh giá
thấp, và kết luận nghiêng về phía tham số. Mục 2.2 của bài Kaplan xác nhận đúng
cách làm ấy, và còn nói thêm một câu: \enquote{results at convergence were
largely independent of learning rate schedule}. Câu ấy nói về điểm hội tụ. Các
mốc giữa chừng thì nó không nói gì, và đó đúng là chỗ Hoffmann lấy dữ liệu ra.

Lý do thứ hai là dải kích thước. Bài viết rằng phần lớn các lần chạy của Kaplan
nhỏ hơn hẳn, nhiều lần dưới 100 triệu tham số
(\enquote{many being less than 100M parameters}), và quan sát rằng đường biên có
\tn{độ cong}{curvature} nhẹ chỉ lộ ra khi đưa các mô hình vài tỷ tham số vào.

\textbf{Và một lý do thứ ba, không có trong bài.} Có một cách kể rằng Kaplan sai
vì trừ embedding ra khỏi \(N\). Tôi đọc cả đoạn giải thích của mục 2 để tìm câu
ấy và nó không có ở đó. Chỗ duy nhất
embedding xuất hiện là phụ lục F, thuần kỹ thuật, nơi bài ghi quy ước của
\emph{chính nó}: \enquote{Note that we also count embeddings matrices in the
total parameter count.} Hai bài đúng là định nghĩa \(N\) ngược nhau, như
mục~\ref{sec:ch09-sau-nd} đã đo, và không bài nào nối chuyện ấy với chỗ số mũ
khác nhau. Một khác biệt định nghĩa có thật, cộng một chỗ bất đồng có thật, ghép
lại thành một lời giải thích mà không ai kiểm.

Ghép hai thứ đúng thành một câu sai là chuyện dễ xảy ra khi cả hai đều tra được.
Đó là lý do mục này trích đoạn văn chứ không tóm tắt nó.
